\documentclass[11pt]{article}
\usepackage[T1]{fontenc}
\usepackage{lmodern,amsmath,amssymb,amsthm,booktabs,microtype}
\usepackage[a4paper,margin=28mm]{geometry}
\usepackage[colorlinks=true,linkcolor=black,citecolor=black,urlcolor=blue,pdfauthor={Jaehun Jung},pdftitle={Interior alternating chains and lower bounds for superpermutations}]{hyperref}
\newtheorem{theorem}{Theorem}
\newtheorem{lemma}[theorem]{Lemma}
\newcommand{\ceil}[1]{\left\lceil #1\right\rceil}
\newcommand{\code}[1]{\texttt{\small #1}}
\setlength{\emergencystretch}{2em}
\title{\bfseries Interior alternating chains and lower bounds\\for superpermutations}
\author{Jaehun Jung\\{\small Independent Researcher, Republic of Korea}}
\date{25 November 2026}
\begin{document}
\maketitle
\begin{abstract}
Let $S(n)$ be the minimum length of a word containing every permutation of an
$n$-letter alphabet as a contiguous substring.
In Pantone's chain-packing bound the return of a pair of alternating chains
controls a linear allowance for their total lengths. We show that the
separator following an interior sequence reduces the permitted length of an
alternating run from $2r$ to $2r-1$ when both orders of the completed pair
are excluded at $r$ repetitions. This gives an explicit bound at least four
letters above Pantone's general formula for every $n\equiv1\pmod 3$, $n\ge16$.
Equality of occupied permutation coordinates supplies earlier exclusions at
$n=10,13,16$. With these exclusions the improvements over the general formula
are respectively $1$, $159$, and $28{,}716$ letters. At every $n\le14$ the existing individual lower bound is at
least as strong as the new general formula. We give a closed formula for all
positive $n$ and a second version using a finite collision test. The word-level inequalities are formalized in Lean.
\end{abstract}

\section{Statement}
Write $P(n)$ for the baseline general bound. For $n\ge8$ it is Pantone's
exact-return formula \cite{pantone,record34}, with its published parameters
and outer integer ceiling. For $1\le n\le7$ we take the values
$1,3,9,33,153,870,5895$. The first two are immediate. The next two
follow from Hunter and Raudvere, the next two from Kisic's inequality
proved by Halfar, and the seventh from Halfar's general bound
\cite{small,halfar}. These choices make it possible to compare a general
formula with stronger proofs specific to one alphabet.

For $n\ge8$ define the chain-packing transfer function
\begin{equation}\label{eq:F}
 F_n(p,t)=n!+(n-1)!+(n-2)!+n-3+
 \ceil{\frac{2(n-2)!-2(n-2)-4p/t}{n^2-4n+1-2(n-1)/t}} .
\end{equation}
On the residue class $n=3m+1$ with $m\ge3$, set $s=3m$.
The public return count in \cite{pantone} is
\[
 R_0=\begin{cases}
  2s/7,&7\mid s,\\
  u(s-u),&7\nmid s,\quad 1\le u<s,\quad 7u\equiv1\pmod s.
 \end{cases}
\]
For any positive $r$ set
\begin{align}
 d_r&=2r-1,&
 p_r&=2m+\tfrac12+\tfrac{1}{2d_r},&
 t_r&=4m-3-\tfrac1{d_r},&
 T_m(r)&=F_{3m+1}(p_r,t_r). \label{eq:prices}
\end{align}
The return count with three explicit collision certificates is
\[
 R_c(m)=
 \begin{cases}3,&m=3,4,\\7,&m=5,\\R_0,&m\ge6.\end{cases}
\]
\begin{theorem}\label{thm:main}
Every word covering every permutation of its $n$ symbols has length at least
\begin{equation}\label{eq:main}
 L(n)=\begin{cases}
 \max\{P(n),T_m(R_0),T_m(R_c(m))\},& n=3m+1,\ m\ge3,\\
 P(n),&\text{otherwise}.
 \end{cases}
\end{equation}
For every $m\ge5$, $L(3m+1)\ge P(3m+1)+4$.
\end{theorem}

The parameters of $P$ outside the improved residue class are those in
\cite{record34}. The exact arithmetic implementation shipped with the paper
includes the small-parameter table and the cycle-order formulas for all
congruence classes. The branch $T_m(R_0)$ already proves the infinite
strict comparison. The branch $T_m(R_c)$ adds the three geometric
certificates in Section~\ref{sec:slots}.

\section{Interior alternation}
We use Pantone's transfer theorem \cite{pantone} in the following form.
In a block-aligned light segment of the permutation overlap path, let
$x_1,\ldots,x_q$ be the lengths of a realized list of chains. An admissible
pair $p,t$ gives \eqref{eq:F} as a word-level lower bound if every such
nonempty list satisfies
\begin{equation}\label{eq:allow}
 2\sum_{i=1}^q x_i\ \le\ 2p+t(q-1).
\end{equation}
The path and its segments here are those extracted from arbitrary covering
words by the inherited packing theorem. In particular, distinct occupied
vertices of the extracted simple path cannot encode the same permutation.
For $n=3m+1$, let $b=2m-1$. A chain has length at most $b$ when $q\ge2$,
and a single chain has length at most $b+1$. These are the pure-return and
basic-return estimates of \cite{pantone}.

In the completed-cell encoding a chain of length $a$ contributes
$C(a)=0^{a-1}1$. The final $1$ is its short separator. In a factor of an
uncompleted realized list, the separator of its last chain may require the
following chain. Call $x_2,\ldots,x_{q-1}$ the interior. A binary alternating
run here is a consecutive list alternating between $b$ and $b-1$.
\begin{lemma}\label{lem:interior}
Suppose each of
$\big(C(b)C(b-1)\big)^r$ and $\big(C(b-1)C(b)\big)^r$
is excluded in every realized segment. Every interior alternating run has
length at most $2r-1$.
\end{lemma}
\begin{proof}
An interior run with at least $2r$ chains has a prefix of $2r$ chains
starting with $b$ or $b-1$. There is a chain after this prefix, even if the
prefix reaches the end of the interior. Hence the separator of the last
chain in the prefix occurs in the realization. Reading the $2r$ completed
chains yields one of the two excluded factors.
\end{proof}

\begin{lemma}\label{lem:allow}
A positive cap $d$ on interior alternating runs gives \eqref{eq:allow}
with $p=b+\frac32+\frac1{2d}$ and $t=2b-1-\frac1d$.
\end{lemma}
\begin{proof}
For $q\ge3$ put $H=\sum_i(b-x_i)$. We recall the two counting inequalities
used in the transfer, writing $B,O,E$ for the ordinary-bad count, singleton
count and excess witness count of Pantone's alternating-chain decomposition
\cite{pantone}. Cutting the interior into binary alternating pieces gives
\[
 q-2\le B+O+d(B+3O+E+1).
\]
Here $B+O$ bounds the removed positions and $B+3O+E+1$ bounds the number
of surviving pieces. The four-chain constraint gives the weighted defect
inequality
\[
 2B+2(b-3)O+E\le 2H-(q-3).
\]
Since $1+d\le2d$ and $1+3d\le2d(b-3)$ for $d\ge1$ and $b\ge5$,
substitution gives $q-2\le d(2H-q+4)$. Using
$\sum_i x_i=bq-H$ yields \eqref{eq:allow}. When $q=2$ the right side is
$4b+2$, whereas $2\sum_i x_i\le4b$. When $q=1$ the pure-return estimate
gives $x_1\le b+1\le p$. The parameter admissibility inequalities of the
transfer hold for $m\ge3$ and $d=2r-1\ge1$ as verified in
the declaration \code{GeneralAlt.tight\_admissible}.
\end{proof}

The return theorem of \cite{pantone} excludes the first orientation at
$r=R_0$. For the second orientation, the two products of the chain
permutations are conjugate, so their $R_0$-th powers are both the identity.
The actual-periodic-pair avoidance lemma therefore excludes the second
orientation as well. Lemmas~\ref{lem:interior} and~\ref{lem:allow} and the
transfer give $T_m(R_0)\le S(3m+1)$.

\section{Earlier collisions}\label{sec:slots}
An occupied slot is specified by $(i,j,k)$, the cell number, row offset
in the cell, and occupied position in that row, each starting at zero.
We normalize the entry permutation of the first cell to the identity.
Later entries are obtained by composing the corresponding inverse cell
permutations. In each row the occupied positions perform prefix rotations
in those normalized coordinates. Two triples in different rows with the
same normalized permutation would therefore realize the same path vertex.
The common starting labeling is a bijection, so this test applies to every
possible entry labeling.

Here is an explicit coordinate recipe for reproducing the test. Use
indices $0,\ldots,s$, and let
\[
 g=(0,s,s-2,s-1,2,1,3,\ldots,s-3),\quad
 h=(0,1,s,2,\ldots,s-1)
\]
be permutations in one-line notation. Starting with $\nu$ equal to the
identity, update $\nu$ to $\nu\circ g$ after a 0-cell and to
$\nu\circ g\circ h$ after a 1-cell. Before this update the zeroth row has
template $(3,4,\ldots,s,2,1,0)$. For positive row offset $j$, rotate the
first $s-1$ symbols of $(1,3,4,\ldots,s-1,2,s,0)$ to the left by $j-1$.
In the resulting template rotate its first $s$ symbols left by $k$ and apply
$\nu$ to all symbols. In a 0-cell use $0\le j\le s-3$ and in a 1-cell
use $j=0,1$. Valid slots are $0\le k<s-2$ for $j=0$, and
$0\le k<s$ for $j>0$. Restricting the 1-cell to two rows suffices for
the witnesses below. Table~\ref{tab:witnesses} gives a pair of equal
occupied permutations for each orientation. The Lean checker compares
every coordinate of each pair.
\begin{table}[ht]
\centering
\begin{tabular}{@{}rrlll@{}}\toprule
$n$ & $r$ & Orientation & Earlier slot & Later slot\\\midrule
10&3&$(b,b-1)$&$(4,0,0)$&$(18,0,1)$\\
  & &$(b-1,b)$&$(8,0,0)$&$(22,0,1)$\\
13&3&$(b,b-1)$&$(0,0,8)$&$(32,0,0)$\\
  & &$(b-1,b)$&$(6,0,8)$&$(38,0,0)$\\
16&7&$(b,b-1)$&$(0,11,11)$&$(109,2,0)$\\
  & &$(b-1,b)$&$(8,11,11)$&$(117,2,0)$\\\bottomrule
\end{tabular}
\caption{Occupied-slot collision certificates in the two completed pair orders.}
\label{tab:witnesses}
\end{table}

These exclusions give $T_m(R_c(m))\le S(3m+1)$ by the two lemmas.
Together with the baseline this proves the first part of Theorem~\ref{thm:main}.
There is also a computable variant. For $1\le r\le R_0$, test the two
completed strings of $r$ pairs for equal occupied permutations in distinct
rows. Set $R_*$ to the first $r$ with a witness in both strings, or $R_0$
when no such pair is found. Replacing $R_c$ in \eqref{eq:main} by $R_*$
again gives a lower bound, by the same proof.

\section{Gain over the general formula}
Put $F=s!$, $c=s(s-2)-2$, and $R=R_0$. Pantone's original slope is
$t_0=4m-3-1/(2R)$ and our slope is $t_1=4m-3-1/(2R-1)$.
In either case $2p+t=8m-2$. Consequently the rational correction before the
ceiling in \eqref{eq:F} is
\[
 Q(t)=\frac{((s-2)t-2)(F-2)-16m}{ct-2s}.
\]
Subtracting the two expressions gives
\[
 Q(t_1)-Q(t_0)=
 \frac{4(F-2)-16mc}
 {(2R)(2R-1)(ct_0-2s)(ct_1-2s)}.
\]
For $s\ge15$ the factors in the denominator are positive. Since
$R\le s^2/4$, $c<s^2$, and $t_i<2s$, their product is less than $s^{10}$.
Also $F\ge s^4\ge8mc+4$, so the numerator is at least $2F$. Finally
$15!\ge2\cdot15^{10}$. For $s\ge15$ one has
$(1+1/s)^{10}\le(5/4)^{10}<16\le s+1$, so induction yields
$s!\ge2s^{10}$. Thus $Q(t_1)-Q(t_0)>2F/s^{10}\ge4$.
Taking the respective ceilings proves the second part of Theorem~\ref{thm:main}.

\section{Evaluations and verification}
Table~\ref{tab:values} compares the closed general formula with its own
baseline and with individual lower bounds recorded in the Superpermutation
Almanac on 25 November 2026 \cite{records}. The relative increases over the general baseline are about $2.48\times10^{-7}$ at ten symbols and $2.36\times10^{-8}$ at thirteen symbols. The two early directions at
thirteen symbols bring the general estimate within 358 letters of the
individual bound. At fourteen the earlier finite-state certificate raises
the lower endpoint a further eight letters \cite{fourteen}.
\begin{table}[ht]
\centering
\small
\begin{tabular}{@{}rrrr@{}}\toprule
$n$&Baseline $P(n)$&Closed bound $L(n)$&Individual lower bound\\\midrule
1&1&1&1\\2&3&3&3\\3&9&9&9\\4&33&33&33\\5&153&153&153\\
6&870&870&872\\7&5,895&5,895&5,899\\8&46,130&46,130&46,134\\
9&408,469&408,469&408,470\\10&4,033,374&4,033,375&4,033,379\\
11&43,917,903&43,917,903&43,917,904\\
12&522,622,380&522,622,380&522,622,401\\
13&6,746,626,531&6,746,626,690&6,746,627,048\\
14&93,891,141,149&93,891,141,149&93,891,141,157\\\bottomrule
\end{tabular}
\caption{Direct evaluations through fourteen symbols. At sizes one and two
the individual values are elementary. Halfar separately reports a
computer-assisted bound of 5,901 at seven symbols \cite{seven}.}
\label{tab:values}
\end{table}

At sixteen symbols the values of $P(16)$ and $L(16)$ are
$22{,}318{,}554{,}295{,}253$ and $22{,}318{,}554{,}323{,}969$.
The companion Lean project links its word-level theorem to the Almanac's
contiguous-substring predicate. The closed formula, its fourteen small
evaluations, the collision test, and the uniform improvement have separate
declarations. The numerical table can also be reproduced with rational
arithmetic in \code{calculation/calculate.py --closed}.

\section*{Parameters of the baseline}
For $n=8,\ldots,12$ use $P(n)=F_n(p,t)$ with $(p,t)$ successively
\[
 (10,46/3),\quad(15,287/12),\quad(521/80,359/40),\quad
 (21,263/12),\quad(28,195/4).
\]
For $n\ge13$ write $n-1=3m+e$ with $e\in\{0,1,2\}$.
For $e=0$ the parameters follow \eqref{eq:prices} with $d=2R_0$.
For $e=1$ use
\[
 p=m(2m+1),\qquad t=2m^2+2m-2-1/R_1.
\]
For $m\bmod6=0,2$ one has $R_1=3m+1$. In the other four classes
$R_1$ is the least common multiple of the indicated two integers.
\[
\begin{array}{c|cccc}
m\bmod6 &1&3&4&5\\ \hline
\text{first}&(2m+1)/3&(7m+3)/6&(2m+1)/3&(7m+1)/6\\
\text{second}&(7m+2)/3&(11m+3)/6&(7m+2)/6&(11m+5)/6
\end{array}
\]
For $e=2$ use $p=(m+1)(2m+1)$ and $t=4m^2+5m-5/2$.
These are the parameters in \cite{pantone,record34}.

\paragraph{Contributions.}
Jaehun Jung proposed transferring the seven-symbol structural investigation
to a general bound and directed the study. OpenAI Codex developed the
additional arguments, generated the certificates, ran the computations
and Lean checks, and drafted the manuscript. Pantone's return and transfer
theorems and the earlier lower bounds are used with attribution.

\begin{thebibliography}{9}
\bibitem{pantone} J. Pantone, \emph{Chain returns and lower bounds for
superpermutations}, proof and manuscript, 2026. \href{https://github.com/jaypantone/2026-10-superpermutation-lower-palomar/tree/ad72ab5010e945b5c8768a918281778c834dcc40}{Pinned source and manuscript}.
\bibitem{record34} Superpermutation Almanac, \emph{Pantone's lower bound},
\href{https://superpermutations.com/b/34}{bound 34}, 2026.
\bibitem{small} Z. Hunter and U. Raudvere, \emph{The Hunter--Raudvere bound},
\href{https://superpermutations.com/b/6}{bound 6}. M. Kisic and J. Halfar,
\emph{Kisic's bound}, \href{https://superpermutations.com/b/7}{bound 7}, 2026.
\bibitem{halfar} J. Halfar, \emph{Halfar's general bound},
\href{https://superpermutations.com/b/9}{bound 9}, 2026.
\bibitem{records} Superpermutation Almanac, \href{https://superpermutations.com/}{individual records},
consulted 25 November 2026.
\bibitem{fourteen} J. Jung, \emph{A finite-state lower bound for superpermutations
on fourteen symbols}, \href{https://superpermutations.com/b/52}{bound 52}, 2026.
\bibitem{seven} J. Halfar, \emph{Seven symbols, lower bound 5,901}, 2026.
\href{https://github.com/jhalfar/superpermutations/tree/8e3c02b5234b49a1422751f86f9e472706082374/proofs/n7_5901}{Proof materials}.
\end{thebibliography}
\end{document}
